\title{LongRoPE: Extending LLM Context Window Beyond 2 Million Tokens}

\begin{abstract}
A sizable context window is an advantageous characteristic for large language models (LLMs). Nevertheless, the existing scope of expanded context windows remains confined to approximately 128k tokens, owing to substantial expenses associated with fine-tuning, limited availability of lengthy textual data, and the detrimental effects that emerge from introducing novel token positions.
\end{abstract}

This work presents LongRoPE, which, for the first time, pushes the context window of pre-trained LLMs to a remarkable \textbf{2048k} tokens. This is accomplished with as few as 1k fine-tuning steps at training lengths up to 256k, all while preserving the model’s capabilities within the standard short context window. The success of LongRoPE relies on three main contributions: (i) we uncover and leverage two distinct types of non-uniformities in positional interpolation through a computationally efficient search procedure, yielding an improved fine-tuning initialization and enabling an 8$\times$ window extension without further fine-tuning; (ii) we propose a stepwise extension approach, wherein we initially fine-tune the model to a 256k context length, and subsequently apply a second round of positional interpolation to the already fine-tuned model to extend the context to 2048k; (iii) we further recalibrate LongRoPE at an 8k context window to restore and maintain the model’s performance on shorter sequences. Comprehensive evaluations conducted on LLaMA2 and Mistral spanning diverse tasks validate the robustness of our approach. Models scaled up via LongRoPE require only slight changes to the positional embedding, preserving the original model architecture and compatibility with most existing optimizations. The implementation will be publicly released at \texttt{https://github.com/microsoft/LongRoPE}

\section{Introduction}

Although Large Language Models (LLMs) have achieved outstanding results across a range of tasks~\cite{gpt4,llama2}, they remain constrained by relatively small context window capacities, such as the 4096-token restriction in LLaMA2~\cite{llama2}. When input lengths exceed this context window, model effectiveness diminishes since the model lacks exposure to these longer positional sequences during training. This limitation creates obstacles in critical applications, including in-context learning involving many examples~\cite{Huang2023FewerIM} and the operation of LLM agents~\cite{park2023generative,madaan2023selfrefine}.

Recent studies have demonstrated that the context window of a pre-trained LLM can be lengthened to approximately 128k through additional fine-tuning with longer documents~\cite{longlora,pi,yarn,zhang2024soaring,liu2023scaling}. Nevertheless, there are three primary challenges hindering further increases to the context window. Firstly, introducing new and previously untrained position indices brings about a multitude of problematic values, resulting in distributional shifts that impede the stability and effectiveness of the fine-tuning process~\cite{pi}. This issue becomes even more pronounced when expanding the window from 4k to over 1000k, as this transition incorporates more than 90\% newly encountered positions. Secondly, successful fine-tuning generally relies on the availability of textual data that matches the intended sequence length. However, existing corpora offer few sufficiently long documents, particularly those surpassing 1000k in length. In addition, the process of training on such ultra-long texts demands substantial computational resources and time, making it impractically expensive in terms of both GPU usage and overall training duration. Thirdly, as context windows are extended to extreme lengths, the model’s attention mechanism becomes increasingly diffuse, being spread over a vast number of token locations, which in turn diminishes the model’s ability to maintain high performance on shorter, original contexts~\cite{pi}.

A common strategy to address the initial challenge involves interpolating RoPE positional embeddings~\cite{rope,pi}, which scales unfamiliar position indices back into the range encountered during pretraining, as illustrated in Fig.\ref{fig:overview}. The Position Interpolation (PI) method~\cite{pi} achieves this by applying linear interpolation to RoPE's rotary angles, guided by the extension ratio. In contrast, NTK~\cite{ntk,dynamicntk} proposes the use of distinct interpolation and extrapolation factors for each RoPE dimension. YaRN~\cite{yarn} further refines this approach by partitioning RoPE dimensions into three frequency-based categories and using extrapolation, NTK, and linear interpolation techniques tailored to each group. Nevertheless, the nature of positional embedding within Transformer architectures is characterized by a \emph{complex, non-uniform distribution of information entropy}. Current methods largely overlook this nuanced variability, resulting in inefficiencies that manifest as information loss and a restricted effective context window.

Section~\ref{sec:analysis} provides empirical evidence for two primary observations: \textbf{(1)} Successful positional interpolation requires addressing two types of non-uniformities: discrepancies in RoPE dimensions as well as disparities among token positions. For lower RoPE dimensions and early token positions, applying less interpolation yields favorable results; however, the best choice is influenced by the desired extension length. \textbf{(2)} By explicitly incorporating these non-uniformities into the positional interpolation process, it becomes possible to better preserve key information inherent in the original RoPE, specifically in crucial dimensions and token positions. This approach reduces information loss due to interpolation and, as a result, produces improved initializations for subsequent fine-tuning. Additionally, this strategy enables extending sequence length up to 8$\times$ without the need for fine-tuning.

Building upon these insights, we present \textbf{LongRoPE}, a robust approach that enables LLMs to operate with context windows surpassing 2 \emph{million} tokens. The design of LongRoPE centers around three core advancements. Firstly, {LongRoPE} leverages the multidimensional nature of non-uniform positional interpolation to its fullest extent. It determines optimal rescale factors for the rotation angles of RoPE, tailored to each dimension and informed by the position of tokens. Since the process of discovering appropriate rescale factors becomes exponentially more complex as the extension ratio increases, {LongRoPE} incorporates an evolutionary search algorithm enhanced with two optimization strategies to accelerate the search process. Fig.~\ref{fig:overview} provides a visualization of the resultant rescaled RoPE obtained through this search.

Subsequently, {LongRoPE} adopts an efficient, staged extension approach that enables a 2048k context window, circumventing the necessity of direct fine-tuning on ultra-long texts—which are rare and typically unavailable. This approach initiates by identifying a 256k context length for the pre-trained LLM, followed by fine-tuning at this specified length. Owing to our non-uniform positional interpolation, which facilitates up to an 8$\times$ extension even without further fine-tuning, we then proceed to determine new RoPE rescale factors on the LLM already fine-tuned for the extended context. Through this method, the desired 2048k context window is successfully established for LLaMA2 and Mistral~\cite{mistral}.

In order to prevent a drop in performance for the original, shorter context windows, {LongRoPE} maintains adjustments to the RoPE rescale factors on the extended LLM. Analogous to the process of increasing from 256k to 2048k, we apply our search algorithm to scale the context window down to 4k and 8k on the 256k fine-tuned LLM, aiming to reduce positional interpolation effects. At inference time, when the input sequence length is below 8k, RoPE is modified using the optimized rescale factors identified during the search.

A comprehensive set of experiments involving multiple LLMs and a range of long-context benchmarks validates the efficacy of our approach. Our findings indicate that LongRoPE consistently preserves low perplexity across evaluation lengths spanning from 4k to 2048k, surpasses 90\% accuracy in passkey retrieval tasks, and achieves accuracy on par with existing methods for standard benchmarks situated within the 4096 context window. Furthermore, LongRoPE is broadly applicable to any LLM utilizing RoPE embedding. The codebase and LongRoPE-2048k models will be made publicly available.

\section{Non-uniformity in Positional Interpolation}
\label{sec:analysis}

\subsection{Preliminary}

Transformer architectures depend on the incorporation of explicit positional cues, typically via position embeddings, to encode the sequence of input tokens. In this study, we center our analysis on the RoPE~\cite{rope} position embedding mechanism, a prevalent method in modern LLMs. For a token located at position index $n$, its associated RoPE encoding can be succinctly described as:

\begin{equation}
		\fontsize{7.8}{7.8} \selectfont
	\label{eq:rope}
	\left[ cos(n\theta_0), sin(n\theta_0), cos(n\theta_1), \cdot\cdot\cdot, cos(n\theta_{d/2-1}), sin(n\theta_{d/2-1}) \right]   
\end{equation}

Here, $d$ denotes the embedding dimension, 
$n\theta_i$ specifies the rotary angle associated with the token at position $n$,
and $\theta_i=\theta^{-2i/d}$ encodes the rotation frequencies. In RoPE, the base value for $\theta$ is typically set to 10000.

\textbf{\emph{Context window extension ratio $s$} and positional interpolation}. We introduce $s$ as the quotient of the extended context window $L'$ to the initial context length $L$: $s=\frac{L'}{L}$.

In order to expand the context window from $L$ to $L'$, existing positional interpolation techniques typically involve reducing the rotation frequencies $\theta_i$ according to the extension ratio $s$. Defining $\beta=\theta^{2/d}$ and letting $\lambda$ represent the specific rescaling factor associated with $s$, we generalize these positional interpolation approaches in the following manner:

\begin{equation}	
	\fontsize{7.5}{7.5} \selectfont
	\label{eq:unified_rope}
	\begin{aligned}
		\left[cos\left(\frac{n}{\lambda(\beta)^0}\right), sin \left(\frac{n}{\lambda(\beta)^0}\right), cos\left(\frac{n}{\lambda(\beta)^1}\right), \cdot\cdot\cdot,  sin \left(\frac{n}{\lambda(\beta)^{d/2-1}}\right) \right]   
	\end{aligned}
\end{equation}

\textbf{Linear positional interpolation (PI)}. PI~\cite{pi} employs a method of linearly interpolating position indices that lie within the original pre-trained context window. When aiming to extend the context length by a ratio $s$, all position rotation angles are uniformly scaled down by a factor of $\lambda=s$ for each RoPE dimension. This uniform compression of positional encoding causes the representations of different positions to become tightly clustered, making it challenging for the model to differentiate between tokens that are near each other in the sequence. As a result, the effectiveness of PI diminishes when applied to substantial context length extensions.

\textbf{NTK-based interpolation and extrapolation}. ~\cite{ntk,dynamicntk} approach RoPE through the lens of information encoding, leveraging Neural Tangent Kernel (NTK) theory~\cite{jacot2018neural,tancik2020fourier} to inform their methods. To address the issue of positional overlap inherent in PI, these works propose spreading the interpolation demand among the different RoPE dimensions. Specifically, they downscale the influence of lower (higher frequency) dimensions while upscaling higher (lower frequency) ones, enabling both positional interpolation and extrapolation via $\lambda=s^{i}$. In the refined dynamic NTK approach~\cite{dynamicntk}, the degree of extension is adapted dynamically at each position depending on the current sequence length. Whereas PI generally requires model fine-tuning, NTK-informed techniques can expand the context window in situations without fine-tuning, though the typical upper bound for extension is around 4$\times$.

\textbf{YaRN}~\cite{yarn} presents a major enhancement in the efficacy of positional interpolation by segmenting the RoPE dimensions into three distinct categories according to their frequency components, with each group subjected to a separate interpolation technique. Dimensions characterized by high frequencies are extrapolated ($\lambda$=1), whereas those corresponding to low frequencies utilize linear interpolation (PI). Intermediate RoPE dimensions make use of the NTK approach. The central innovation in YaRN is this method of organizing RoPE dimensions; however, the grouping process currently relies on manual, empirically-driven heuristics, which may lead to less than optimal outcomes when applied to new LLM architectures.

\subsection{Study on Non-uniform Positional Interpolation}

Drawing upon insights from NTK and YaRN, it becomes apparent that their improvements are largely attributed to leveraging non-linearity—particularly through managing distinct frequencies along RoPE dimensions for targeted interpolation and extrapolation. Yet, the prevailing non-linear mechanisms are largely dictated by manually crafted heuristics. This observation prompts two key inquiries: (1) Does the present method for positional interpolation represent the optimal approach? (2) Might there exist additional non-linearities that have not yet been investigated?

In order to address these questions, we employ evolution search (refer to Sec.~\ref{sec:method}) to identify improved non-uniform positional interpolation schemes for LLaMA2-7B. The process is steered by perplexity measurements, utilizing 5 randomly selected examples from the PG19~\cite{pg19} validation dataset. Our experimental results yield several important insights.

\textbf{Finding 1}: \textit{RoPE dimensions exhibit substantial non-uniformities, which are not effectively handled by current positional interpolation methods.}

We conduct a search for the optimal $\lambda$ value for each RoPE dimension as described in Eq.~\ref{eq:unified_rope}. Table~\ref{tbl:exp1} presents a comparison of the perplexity values for LLaMA2-7B across various methods on the PG19 and Proof-pile~\cite{proof-pile} test sets, evaluated without fine-tuning. The approach identified through our search yields notable gains, indicating that both the linear (PI) and non-uniform (Dynamic-NTK and YaRN) interpolation strategies employed previously are not optimal. Interestingly, on the PG19 dataset, YaRN performs worse than both PI and NTK, likely because it does not attain the designated context window length in the absence of fine-tuning for the LLM. Specifically, in the case of an 8k context size, the perplexity with YaRN increases sharply after reaching 7k tokens.

In our investigation, we observe that the rescaled factors $\lambda$ in Eq.~\ref{eq:unified_rope} are not uniform, in contrast to the constant scale $s$ employed in PI, the formula used in NTK, and the group-based scaling of YaRN. This non-uniform rescaling yields a notable enhancement in the language modeling capabilities (measured by perplexity) of LLaMA2 for 8k and 16k context lengths, achieved without additional fine-tuning. The improvement arises because these positional embeddings retain much of the original RoPE structure—particularly along critical dimensions—which facilitates the model's ability to differentiate between nearby token positions.

\textbf{Finding 2}: \textit{RoPE for the initial tokens in the input sequence should be extrapolated with less interpolation.} 

We propose that for the first $\hat{n}$ tokens within input sequences, the RoPE mechanism should employ minimal interpolation. The rationale is that these tokens are associated with substantial attention weights, rendering them especially significant within attention layers—a finding supported by prior work such as Streaming LLM~\cite{streamingllm} and LM-Infinite~\cite{han2023lminfinite}. To empirically investigate this, we expand the context length to 8k and 16k tokens using both PI and NTK, systematically varying the count of initial tokens ($\hat n$ = 0, 2, ..., 256) that are exempted from interpolation. For $\hat n = 0$, the method defaults to standard PI and NTK configurations. The experimental outcomes, summarized in Table~\ref{tbl:tokenposition}, yield two primary conclusions: \textbf{(1)} preserving the initial tokens without interpolation of positional information leads to notable performance gains, and \textbf{(2)} the ideal selection of $\hat n$ is contingent upon the desired extension length of the context.

\textbf{Finding 3}: \textit{Non-uniform positional interpolation effectively extends LLM context window in both fine-tuning and non-fine-tuning settings.} 

Our experiments demonstrate that the non-uniform position interpolation derived from our search process notably boosts extension capabilities at 8k and 16k token lengths without additional fine-tuning; however, achieving even greater extension to longer contexts necessitates fine-tuning. Accordingly, we fine-tune LLaMA2-7B with the optimized RoPE found through our search for a context window of 64k tokens (see Appendix for the detailed configuration). Table~\ref{tbl:finetune} illustrates that our approach consistently outperforms both PI and YaRN, regardless of whether LLaMA2-7B is fine-tuned. This advantage stems from our method's strategic non-uniform positional interpolation, which effectively preserves information and provides a superior initialization for the subsequent fine-tuning process.

\textbf{Summary}. Our findings identify two forms of non-uniformity: variability among RoPE dimensions and across token positions. Exploiting these types of non-uniformity in positional interpolation yields substantial gains in LLM context length extension.

\section{LongRoPE}

\label{sec:method}
Building upon these insights, we propose LongRoPE, which initially incorporates a search algorithm designed to efficiently leverage both types of non-uniformity. Utilizing this approach, LongRoPE successfully extends the LLM context window to exceed 2 million tokens.

\subsection{Problem Formulation}

These two forms of non-uniformity expand the potential solution landscape, thereby complicating the optimization process. To mitigate these challenges, we recast the multidimensional non-uniform position interpolation optimization as a search-based problem.

Consider an LLM designed for a context window of size $L'$ and input documents $\mathbf{X}$ where every sequence $\mathbf{x} \in \mathbf{X}$ has a token length exceeding $L'$. Let the initial rotary angle for the $i^{th}$ dimension within RoPE embeddings at the $n^{th}$ token position be expressed as $\frac{n}{\beta^i}$. Under this setting, the optimization task is articulated as follows:

\begin{equation}
\begin{gathered}
		\label{eq:problem}

\underset {\mathbf{x}\in\mathbf{X}; \, |\mathbf{x}|\ge L'}{ \operatorname {arg\,min} } \,  \mathcal{L}\, \left( \text{LLM} (\text{RoPE}, \mathbf{X})\right),	\text{with \;} 
\\
\scalebox{0.93}{
$\underset {\begin{subarray}{l} i=0,\cdot\cdot,\frac{d}{2}-1;\\ n\in[ 0,|\mathbf{x}|);\end{subarray}}{\text{RoPE}_(n)}= {\left[\cdots, \cos\left(\mathbb{I}(\hat{\lambda}_{i}, \hat{n})\frac{n}{\beta^i}\right), \sin\left(\mathbb{I}(\hat{\lambda}_{i}, \hat{n})\frac{n}{\beta^i}\right), \cdots\right] }$}\\
	\text{where \;} \mathbb{I}(\hat{\lambda}_{i},\hat{n})=\begin{cases}
	1&\text{ if } n < \hat{n}\\
{\frac{1}{\lambda}_{i}}&\text{ if } n \geq \hat{n}
\end{cases}
	\end{gathered}
\end{equation}

In this approach, we define a collection of scaling coefficients, $\mathbb{I}(\hat{\lambda}_{i},\hat{n})$, designed to address both types of non-uniformities present. The variables $\hat{\lambda}_{i}$ and $\hat{n}$ represent the non-uniformity in RoPE dimensions and in the positions of tokens, respectively. More precisely, $\mathbb{I}(\hat{\lambda}_{i},\hat{n})$ functions to adjust the rotation angle for the $i^{th}$ RoPE dimension; here, $\hat\lambda_{i}$ serves as the scaling coefficient, while $\hat{n}$ marks the positional threshold. For token positions before $\hat{n}$ (that is, for the first $\hat{n}-1$ tokens), the scaling factor $\hat\lambda_{i}$ is not employed, maintaining the standard RoPE rotation angle $\frac{n}{\beta^i}$. In contrast, when token positions reach or exceed $n\ge\hat{n}$, the scaling coefficient is introduced.

Assuming a desired context window length of $L'$, our task is to determine the most suitable set of rescale coefficients ($\mathbb{I}(\hat{\lambda}_{0},\hat{n})$, $\mathbb{I}(\hat{\lambda}_{1},\hat{n})$, ..., $\mathbb{I}(\hat{\lambda}_{i},\hat{n})$, ...) for each RoPE dimension from $1^{st}$ to $d^{th}$. By applying these optimized rescale factors to the $\text{RoPE}$ of the target $\text{LLM}$, the model is able to minimize the next-token prediction loss, denoted by $\mathcal{L}$ (perplexity), when processing input sequences $\mathbf{X}$ that have a length of $L'$.

\subsection{Searching the Non-uniform Position Interpolation}

To address the challenge posed in Eq.~\ref{eq:problem}, we present a straightforward yet powerful approach that identifies the optimal RoPE rescale factors to make the most of the multidimensional disparities in position embedding.

\textbf{Search space}. Our approach encompasses a broad search space to capture both forms of non-uniformity, as outlined in Table~\ref{tbl:searchspace}. In particular, we permit the independent optimization of rescale factors for each RoPE embedding dimension. To facilitate a more tractable search space, our method searches over $\lambda_i$ and $\hat{n}$, as opposed to directly searching for $\mathbb{I}(\hat{\lambda}_{i},\hat{n})$, with the relationship $\hat{\lambda}_{i}=1/\lambda_i$. According to Table~\ref{tbl:searchspace}, the parameter $\lambda_i$ may take on values ranging from 1.0 (equivalent to basic extrapolation) up to $s\times1.25$ (corresponding to broader interpolation than that of PI), incremented by 0.01, where $s$ denotes the desired extension ratio for the context window.

The parameter $\hat{n}$ determines how many of the earliest token positions preserve their original RoPE embedding, foregoing position interpolation. In practice, we let $\hat{n}$ vary over the set \{0, 1, 2, 4, 8, 12, 16, 20, 24, 28, 32, 64, 128, 256\} during the search. If $\hat{n}=0$, then every token position employs the rescale factors obtained through the search.

\textbf{Evolution-based search}. The search space detailed in Table~\ref{tbl:searchspace} encompasses a vast array of positional interpolation options, making exhaustive exploration highly impractical. For instance, when considering an extension of $s=4\times$, the number of possible combinations reaches $400^{128/2}\times14=4\times10^{167}$. As the extension ratio increases, the size of the search space grows exponentially. To effectively navigate this immense space, we employ evolution search~\cite{guo2020single} and incorporate two optimization strategies that significantly improve search efficiency. The overall methodology for the search process is outlined in Algorithm~\ref{alg:search}.

\textit{Optimized initial population generation}. Rather than starting with a purely random population of $P$ rescale factors, we explicitly include the three RoPE rescale factors representing PI, NTK, and YaRN among the initial individuals. The other $P$-3 members of the population are produced by randomly applying mutations to these three rescale factors with probability $p$.

\textit{Monotonically non-decreasing constraint}. Once the initial population has been produced, LLM perplexity is evaluated for every individual. This involves adjusting the target LLM using the designated RoPE rescale factors and determining the perplexity on the input $\mathbf{X}$. The best k individuals, ranked by perplexity, are selected to serve as parents in the evolutionary process. Yet, due to the expansive nature of the search space, simple application of mutation and crossover often results in poor candidate solutions, which in turn necessitate superfluous perplexity evaluations. This issue is exacerbated as $L'$ grows, since each perplexity assessment requires considerable computational resources.

To mitigate this issue, we enforce a constraint ensuring that the sampled RoPE rescaling factors are monotonically non-decreasing, specifically $\lambda_i\le \lambda_{i+1}$. Only those RoPE parameterizations that fulfill this monotonicity condition are used during perplexity assessments with the LLM, which markedly narrows the search space and enhances efficiency. More precisely, we mandate that $\lambda_i$ progresses in a non-decreasing manner across the RoPE dimensions, indexed by $i = 0, \ldots, 63$. This requirement stems from insights offered by NTK theory~\cite{jacot2018neural,tancik2020fourier,ntk}, which indicate that dimensions associated with higher frequencies (lower $i$) should have reduced interpolation (i.e., smaller values of $\lambda_i$), whereas dimensions tied to lower frequencies (higher $i$) are suited for increased interpolation (i.e., larger $\lambda_i$).

\begin{algorithm}[t]
	\small
	\caption{The search algorithm}
	\textbf{Input:} target LLM, input samples $\mathbf{X}$, population size $P$, mutation size $N_1$, crossover size $N_2$, maximum iterations $\mathcal{T}$, mutation probability $p$\\

	\begin{algorithmic}[1]
		\label{alg:search}
		\STATE Top-k $\gets \phi$;	
		\STATE $\text{P}_0 \gets$ \textit{Initialize\_population\_with\_optimization} ($P$, $\mathbf{X}$, $p$); 
		\FOR{$i = 1$ to $\mathcal{T}$}
		\STATE \textit{Compute\_perplexity} (LLM, $\text{P}_{i-1}$, $\mathbf{X}$);
		\STATE Top-k $\gets$ \textit{Update\_Topk}(Top-k, $\text{P}_{i-1}$);
		\STATE $\text{P}_{mutation} \gets$ \textit{Mutation\_with\_mono\_constraint}(Top-k, $N_1$, $p$); 
		\STATE $\text{P}_{crossover} \gets$ \textit{Crossover\_with\_mono\_constraint}(Top-k, $N_2$);
		\STATE $\text{P}_i \gets \text{P}_{mutation} \cup \text{P}_{crossover} \cup$ Top-k;
		\ENDFOR
		\STATE Output the individual with the minimal perplexity among Top-k;
	\end{algorithmic}
\end{algorithm}

\textbf{8$\times$ extension without fine-tuning}. Through evolutionary search, we efficiently discover non-uniform RoPE rescale factors that retain important dimensions and positions, thereby reducing the information loss typically associated with interpolation. As shown in Fig.\ref{fig:non-ft}, our approach successfully expands LLaMA2's context window from 4k to 32k tokens without the need for any fine-tuning. In comparison, current techniques like PI, and non-uniform NTK and YaRN exhibit a sharp increase in perplexity after only a 2$\times$ extension.

\subsection{Extending LLM Context Window to 2048K}
\label{sec:extendto2048k}

\textbf{Progressive extension to 2048k}. In this section, we present our approach for expanding the context window of pre-trained LLMs from the standard 4k up to beyond 2048k tokens. As shown, by employing our non-uniform positional interpolation, an extension by a factor of 8$\times$ is attainable without requiring any fine-tuning. For even larger increases, such as a 512$\times$ extension, fine-tuning becomes indispensable. A common strategy involves selecting appropriate RoPE rescaling factors tailored to the desired 2048k context length, followed by fine-tuning. Nevertheless, this approach is hindered by the substantial computational costs it entails. Furthermore, our empirical observations indicate that fine-tuning LLMs at such extreme extension ratios is notably difficult (see Appendix).

LongRoPE demonstrates robust performance with both the baseline and further trained extended LLMs. Accordingly, we propose a streamlined, incremental approach that attains the desired 2048k context length in only 1k fine-tuning iterations, while restricting training to sequences of up to 256k tokens.

$\Diamond$ \textit{Exploration of pre-trained LLM adaptation to 256k context length using LongRoPE search}. Using LLaMA2 as a case study, we systematically examine search procedures for desired context window lengths of 128k and 256k. At this phase, the expansion ratios achieved are 32$\times$ and 64$\times$, respectively.

$\Diamond$ \textit{Fine-tuning to 256k}. Next, we adapt the pre-trained LLM to support a 256k context window. Our approach begins by fine-tuning LLaMA2 for 400 iterations utilizing RoPE rescaling parameters set for 128k. After this phase, we switch the RoPE rescaling to correspond to 256k on the resulting checkpoint, followed by a further 600 fine-tuning steps. Compared to directly fine-tuning for 256k, this staged process demonstrates greater efficiency.

$\Diamond$ \textit{Scaling a fine-tuned extended LLM up to 2048k via LongRoPE search}. In the last step, we conduct an additional search on the LLM that was previously fine-tuned with a 256k context length. This process successfully expands the context window to an unprecedented 2048k, all without the need for extra fine-tuning. The resulting context window extension represents a $512\times$ increase.

\textbf{Original context window performance degradation}.  When expanding the context window to a substantial 2048k, a decline in model performance is observed within the initial context window. This phenomenon has been documented in the literature on positional interpolation~\cite{pi}, since projecting position embeddings into significantly higher-dimensional space results in their compression into a compact interval for the original context range, thereby impairing model effectiveness. Employing a 512$\times$ context window extension exacerbates this effect, causing embeddings within the standard 4k window to overlap considerably.

To address this issue, we conduct an additional evolutionary search on the augmented LLM, fine-tuning the RoPE rescale parameters for shorter context lengths such as 4k and 8k. Given that these shorter contexts require less positional interpolation, we lower the upper bound for the $\lambda$ values explored in the search. At inference time, the LLM adapts the RoPE rescale factors on the fly as appropriate for the given context length.

\section{LongRoPE}

\label{sec:method}
Building on these observations, we propose LongRoPE. This approach begins by presenting a search algorithm specifically designed to leverage the two observed non-uniformities efficiently. Subsequently, this algorithm is utilized to expand the context window of LLMs to exceed 2 million tokens.

\subsection{Problem Formulation}

The aforementioned non-uniformities greatly expand the potential solution space and add layers of difficulty to the optimization process. To tackle these challenges, we recast the multidimensional non-uniform position interpolation optimization task as a search problem.

Consider an LLM designed for a context window size of $L'$, handling long input documents $\mathbf{X}$ where each sequence $\mathbf{x}\in\mathbf{X}$ contains more tokens than $L'$. Let the baseline rotary angle at position $n$ for the $i^{th}$ dimension in the RoPE embedding be represented by $\frac{n}{\beta^i}$. We define the optimization task accordingly:

\begin{equation}
\begin{gathered}
		\label{eq:problem}

\underset {\mathbf{x}\in\mathbf{X}; \, |\mathbf{x}|\ge L'}{ \operatorname {arg\,min} } \,  \mathcal{L}\, \left( \text{LLM} (\text{RoPE}, \mathbf{X})\right), \text{where} \;
\\
\scalebox{0.93}{
$\underset {\begin{subarray}{l} i=0,\ldots,\frac{d}{2}-1;\\ n\in[0,|\mathbf{x}|);\end{subarray}}{\text{RoPE}_(n)}= {\left[\ldots,\cos\left(\mathbb{I}(\hat{\lambda}_{i}, \hat{n})\cdot\frac{n}{\beta^i}\right), \sin\left(\mathbb{I}(\hat{\lambda}_{i}, \hat{n})\cdot\frac{n}{\beta^i}\right),\ldots\right] }$}\\
	\text{where} \; \mathbb{I}(\hat{\lambda}_{i},\hat{n})=\begin{cases}
	1&\text{if } n<\hat{n}\\
	{\frac{1}{\lambda}_{i}}&\text{if } n\geq\hat{n}
\end{cases}
	\end{gathered}
\end{equation}

In this approach, we define a group of scaling factors, $\mathbb{I}(\hat{\lambda}_{i},\hat{n})$, to address both categories of non-uniformity present. The variables $\hat{\lambda_i}$ and $\hat{n}$ correspond to the unevenness within RoPE dimensions and token positions, respectively. In practice, $\mathbb{I}(\hat{\lambda}_{i},\hat{n})$ serves to adjust the rotation angle specific to the $i^{th}$ RoPE dimension; here, $\hat\lambda_{i}$ functions as the scaling parameter and $\hat{n}$ designates the positional cutoff. For token positions prior to $\hat{n}$ (i.e., positions less than $\hat{n}$), the scaling factor $\hat\lambda_{i}$ is not employed, retaining the standard RoPE rotary angle $\frac{n}{\beta^i}$. Once token positions reach or exceed $n \ge \hat{n}$, the adjustment via the rescale factor is activated.

Assuming a desired context window length of $L'$, the task is to determine the best set of rescale factors ($\mathbb{I}(\hat{\lambda}_{0},\hat{n})$, $\mathbb{I}(\hat{\lambda}_{1},\hat{n})$, ... ,$\mathbb{I}(\hat{\lambda}_{i},\hat{n})$, ...) across all RoPE dimensions from $1$ to $d$. With these rescaled RoPE parameters, the target $\text{LLM}$ is expected to minimize the next-token prediction loss $\mathcal{L}$ (or equivalently, the perplexity) on input sequences $\mathbf{X}$ of token length $L'$.

\subsection{Searching the Non-uniform Position Interpolation}

To address the problem defined in Eq.~\ref{eq:problem}, we propose a straightforward yet highly efficient technique that identifies optimal RoPE rescale factors, aiming to take full advantage of the multidimensional non-uniformities inherent in position embeddings.

\textbf{Search space}. We construct an extensive search space that captures both forms of non-uniformity. An overview of this search space is provided in Table~\ref{tbl:searchspace}. In particular, our approach entails learning an individualized rescale factor for each RoPE embedding dimension. To simplify the structure of the search space, we opt to search over $\lambda_i$ and $\hat{n}$, rather than searching directly for $\mathbb{I}(\hat{\lambda}_{i},\hat{n})$; here, $\hat{\lambda}_{i}=1/\lambda_i$. As detailed in Table~\ref{tbl:searchspace}, the search interval for $\lambda_i$ ranges from a lower bound of 1.0 (corresponding to standard extrapolation) up to an upper limit of $s\times1.25$ (permitting interpolation beyond that used in PI), using increments of 0.01, where $s$ denotes the ratio of the target extension in context window length.

The parameter $\hat{n}$ determines how many of the initial token positions preserve their original RoPE embedding, foregoing position interpolation. In our experiments, $\hat{n}$ is permitted to vary among the set \{0, 1, 2, 4, 8, 12, 16, 20, 24, 28, 32, 64, 128, 256\}. Notably, if $\hat{n}=0$, then every token position relies on the optimized rescale factors rather than the original embedding.

\textbf{Evolution-based search}. The search space outlined in Table~\ref{tbl:searchspace} encompasses a vast array of positional interpolation options, making efficient navigation particularly demanding. For instance, when considering an extension factor of $s=4\times$, the total number of possible configurations reaches $400^{128/2}\times14 = 4\times10^{167}$. As the extension ratio increases, the search space grows at an exponential rate. To efficiently manage this complexity, we employ evolution search as proposed in~\cite{guo2020single} and integrate two optimization strategies that significantly enhance the efficiency of the search process. The overall search methodology is detailed in Algorithm~\ref{alg:search}.

\textit{Optimized initial population generation}. In place of a fully random initialization for the population of $P$ rescale factors, we explicitly include the RoPE rescale factors associated with PI, NTK, and YaRN as specific members of the starting population. The rest of the $P-3$ individuals are formed by introducing random mutations to these three rescale factors, applying mutations with probability $p$.

\textit{Monotonically non-decreasing constraint}. Once the initial population has been established, we evaluate the LLM perplexity of every individual. This process involves applying each candidate’s designated RoPE rescale factors to the target LLM and measuring the perplexity on the input $\mathbf{X}$. The highest-ranked $k$ individuals are then selected to serve as the parent pool for subsequent evolutionary steps. Nevertheless, because of the immense size of the search space, indiscriminate mutation and crossover operations may result in the exploration of suboptimal candidates, thereby triggering excessive and avoidable perplexity evaluations. This inefficiency is amplified for large $L'$, as computing perplexity for each candidate entails a substantial inference cost.

To tackle this issue, we enforce a monotonicity constraint such that the rescaled RoPE factors sampled must satisfy $\lambda_i\le \lambda_{i+1}$. Only those RoPE configurations that adhere to this requirement are considered during LLM perplexity assessment, thus markedly narrowing the search space. More precisely, our approach mandates that $\lambda_i$ forms a non-decreasing sequence across the RoPE dimensions (with $i=0,...,63$). The rationale for this dimensional monotonicity stems from NTK theory~\cite{jacot2018neural,tancik2020fourier,ntk}, which indicates that higher-frequency components, associated with lower dimensions, benefit from reduced interpolation (smaller $\lambda_i$), whereas lower-frequency, higher-dimensional components are amenable to greater interpolation (larger $\lambda_i$).

\begin{algorithm}[t]
	\small
	\caption{The search algorithm}
	\textbf{Input:} target LLM, input samples $\mathbf{X}$, population size $P$, mutation size $N_1$, crossover size $N_2$, max iterations $\mathcal{T}$, mutate probability $p$\\

	\begin{algorithmic}[1]
		\label{alg:search}
		\STATE Top-k $\leftarrow \phi$;
		\STATE $\text{P}_0$ $\leftarrow$ \textit{Initialize\_population\_with\_optimization}($P$, $\mathbf{X}$, $p$);
		\FOR{$i = 1$ \textbf{to} $\mathcal{T}$}
		\STATE \textit{Compute\_perplexity}(LLM, $\text{P}_{i-1}$, $\mathbf{X}$);
		\STATE Top-k $\leftarrow$ \textit{Update\_Topk}(Top-k, $\text{P}_{i-1}$);
		\STATE $\text{P}_{mutation}$ $\leftarrow$ \textit{Mutation\_with\_mono\_constraint}(Top-k, $N_1$, $p$);
		\STATE $\text{P}_{crossover}$ $\leftarrow$ \textit{Crossover\_with\_mono\_constraint}(Top-k, $N_2$);
		\STATE $\text{P}_i$ $\leftarrow$ $\text{P}_{mutation}$ $\cup$ $\text{P}_{crossover}$ $\cup$ Top-k;
		\ENDFOR
		\STATE Return the member of Top-k exhibiting the minimum perplexity;
	\end{algorithmic}
\end{algorithm}

\textbf{8$\times$ extension without fine-tuning}. Through evolutionary search, our approach successfully discovers non-uniform RoPE rescale factors that maintain crucial positional and dimensional features, thereby reducing information loss from interpolation. As shown in Fig.\ref{fig:non-ft}, this technique enables LLaMA2's context window to be expanded from 4k up to 32k tokens without necessitating fine-tuning. By comparison, other approaches like PI, as well as non-uniform NTK and YaRN, result in a sharp increase in perplexity beyond a 2$\times$ extension.

\subsection{Extending LLM Context Window to 2048K}
\label{sec:extendto2048k}

\textbf{Progressive extension to 2048k}. In this section, we present our approach for scaling the context window of pre-trained LLMs from the standard 4k length to beyond 2048k. Our results indicate that employing non-uniform positional interpolation allows for an 8$\times$ expansion without requiring any fine-tuning. When seeking to achieve substantially larger extensions (such as 512$\times$), however, fine-tuning becomes indispensable. A commonly considered strategy involves optimizing RoPE rescaling factors for the desired 2048k window, followed by a fine-tuning phase. This approach, however, is often hindered by the substantial computational costs involved. Furthermore, according to our empirical observations, effectively fine-tuning LLMs for such high extension ratios presents significant difficulties (refer to Appendix).

Fortunately, LongRoPE demonstrates efficacy with both the base and fine-tuned versions of the extended LLM. As a result, we propose a progressive and efficient strategy that reaches the desired 2048k context window in only 1k fine-tuning iterations, using sequences up to 256k tokens during training.

$\Diamond$ \textit{LongRoPE search for scaling pre-trained LLMs to 256k}. Using LLaMA2 as a case study, we systematically explore configuration settings targeting context window lengths of 128k and 256k. At this phase, the model undergoes context extension by factors of 32$\times$ and 64$\times$, accordingly.

$\Diamond$ \textit{Fine-tuning to 256k}. Next, the pre-trained LLM undergoes further fine-tuning to support a context window length of 256k. In detail, we initially perform 400 fine-tuning steps on LLaMA2 utilizing RoPE rescaling parameters set for 128k. Upon completion, we switch the RoPE rescaling parameters to those corresponding to 256k on the obtained checkpoint and continue with an extra 600 fine-tuning steps. This staged procedure demonstrates greater efficiency compared to attempting to fine-tune to 256k in a single process.

$\Diamond$ \textit{Scaling fine-tuned extended LLM to 2048k using LongRoPE search}. In the last step, we apply a subsequent search procedure to the LLM that was previously fine-tuned for 256k sequence lengths. This process enables the context window to expand to a massive 2048k tokens, achieved without any additional rounds of fine-tuning. As a result, the total context length is extended by a factor of $512\times$.

\textbf{Recovery within shorter context windows}. Upon increasing the context window to a very large 2048k tokens, we observe a decline in performance on tasks using the standard, shorter context window. This reduction in effectiveness is a recognized limitation of positional interpolation~\cite{pi}, since it compresses position embeddings for the original window into a more confined subspace, thereby impairing the language model's accuracy. When scaling up by a factor of 512, the positional representations within the initial 4k context window are especially densely packed.

To address this issue, we conduct an additional evolution search on the expanded LLM to fine-tune RoPE rescale factors tailored to shorter context lengths (such as 4k and 8k). For these reduced lengths, we constrain the maximum permissible $\lambda$ in the search process, as minimal positional interpolation is necessary. At inference time, the LLM adaptively modifies the relevant RoPE rescale factors accordingly.

 \section{Experiments}

\label{sec:exp}
\subsection{Setup}
\textbf{Evaluation Tasks and models}. We deploy LongRoPE with both LLaMA2-7B and Mistral-7B, and assess its capabilities across three dimensions: (1) measuring the perplexity of large language models with extended context on lengthy documents; (2) the Passkey retrieval task, designed to test the model’s proficiency in extracting a straightforward passkey amidst substantial irrelevant content; and (3) performance on conventional LLM benchmark tasks confined to a 4096 token context window.

\textbf{Fine-tuning}. In the case of LLaMA2, we adopt a learning rate of 2e-5 that decays linearly, and set the global batch size to 32. The initial fine-tuning phase comprises 400 steps using the Redpajama~\cite{redpajama} dataset, segmented into 128k token blocks that are framed by the BOS and EOS tokens. Subsequently, starting from the resulting checkpoint, we perform an additional 600 training steps to extend the context length to 256k. Training with a 128k context window utilizes 8 A100 GPUs, facilitated by the distributed training system~\cite{cube}, whereas scaling up to 256k context requires 16 A100 GPUs. For Mistral, we use a fixed learning rate of 1e-6 with a global batch size of 64. Both the 128k and 256k context models are trained in alignment with the YaRN~\cite{yarn} methodology, running for 400 steps on Together Computer's Long-Data Collections~\cite{mistraldata} at a 16k sequence length, and leveraging 4 A100 GPUs for computation.

\textbf{Search}. When the target window size is up to 256k, the following configuration is employed: $P$=64, $N_1$=$N_2$=16, $p$=0.3, $\mathcal{T}$=40, with the top-32 candidates from each iteration advancing for mutation and crossover. Perplexity evaluation utilizes 5 randomly chosen samples from the PG19 validation set, each meeting or exceeding the target context length. For windows exceeding 512k, the numbers for population, mutation, and crossover are reduced by half. In this scenario, perplexity is assessed on 3 random validation samples from Pile-Books3~\cite{pile}.

\textbf{Baselines}. To achieve a 2048k context window, we fine-tuned our models using context lengths of 128k and 256k. As a result, we obtained LongRoPE-2048k (ft=128k) for LLaMA2 and LongRoPE-2048k (ft=256k) for Mistral. For evaluation, we benchmark these models against leading approaches for context window extension. Specifically, we include open-source large language models that have undergone fine-tuning after positional interpolation via PI, NTK, and YaRN methods. The baseline models for comparison consist of Together-32k~\cite{together}, Code LLaMA~\cite{codellama}, LongLoRA-full-FT-100k~\cite{longlora}, as well as YaRN-LLaMA and YaRN-Mistral~\cite{yarn}.

\subsection{Main Results}

\label{sec:mainresult}
\textbf{Language modeling of long sequences up to 256k}. Our analysis begins with a comparison against the most advanced long-context LLMs at an evaluation length of 256k tokens. To establish the robustness of our approach, we utilize two datasets: the Proof-pile~\cite{pg19} and PG19~\cite{pile} test splits. Perplexity is assessed across different context lengths using a 256-token sliding window. The entirety of the PG19 test split, which comprises 100 documents, is employed. In the case of Proof-pile, following the approach in YaRN~\cite{yarn}, we randomly choose 10 samples, each having a minimum length of 128k tokens.

Table~\ref{tbl:proof-pile} and Table~\ref{tbl:pg19} present a perplexity comparison for LLaMA2 and Mistral models augmented using various interpolation approaches, evaluated on the Proof-pile and PG19 datasets, respectively. Two primary findings emerge from these results: \textbf{(1)} the perplexity of our extended models consistently declines as the evaluation length increases from 4k to 256k, demonstrating their capability to effectively utilize extended context. \textbf{(2)} Notably, even with a context window that is 16$\times$ larger—a scenario often associated with reduced performance at shorter contexts—our LongRoPE-2048k models achieve superior results compared to leading baselines within a 256k context window.

\textbf{Language modeling for sequences exceeding 2000k}. To assess performance on very long texts, we utilize the Books3~\cite{pile} corpus. For the sake of efficient evaluation, we randomly sample 20 books with lengths greater than 2048k and apply a 256k sliding window approach.

As illustrated in Table~\ref{tbl:books3}, LongRoPE is capable of extending the context window of both LLaMA2-7B and Mistral-7B to 2048k, while maintaining perplexity levels on par with, or surpassing, those of the baseline models for shorter context lengths ranging from 8k to 128k. There are clear disparities in performance between the 2048k versions of LLaMA2 and Mistral. Mistral demonstrates superior results compared to baselines at shorter window lengths, yet its perplexity climbs above 7 for contexts beyond 256k. The results for LLaMA2 are as anticipated: perplexity consistently declines with increasing context lengths, although it shows slight rises at 1024k and 2048k. Additionally, on LLaMA2, fine-tuning LongRoPE-2048k at 256k yields improved performance relative to fine-tuning at 128k, attributable to the reduced secondary extension ratio (8$\times$ as opposed to 16$\times$). Conversely, Mistral achieves better outcomes when the fine-tuning context size is set to 128k. This can be largely explained by the training protocol, as we adopt the YaRN configuration of a 16k training length for Mistral's 128k and 256k fine-tuning, which in turn constrains Mistral’s capability to extend its context window in subsequent stages.

\textbf{Passkey retrieval}. Next, we investigate the practical context window length in generation scenarios. We adopt the synthetic passkey retrieval task introduced by ~\cite{passkey}. In this setup, the model must locate a randomly chosen passkey—a sequence of five digits—embedded somewhere within an extensive document. Further information on the prompt format can be found in the appendix. We conduct 10 trials of this retrieval exercise, each time inserting the passkey at a uniformly random position throughout the evaluation context length.

Fig.~\ref{fig:passkey} presents a comparison of retrieval accuracy against baseline models. While the accuracy of previous approaches plummets to zero for sequence lengths above 128k, our LongRoPE-LLaMA2-2048k (ft=256k) achieves consistently high retrieval accuracy ($\ge$90\%) across an extensive range from 4k up to 2048k, even when tasked with retrieving a passkey from millions of tokens. Similarly, LongRoPE-Mistral-2048k (ft=128k) maintains perfect accuracy up to 1800k, only decreasing to 60\% at 2048k—a trend that mirrors the results reported in Table~\ref{tbl:books3}, where a modest increase in perplexity is observed at the 2048k context window.

\textbf{Evaluation on established benchmarks within the original context window}. We assess the performance of LongRoPE-2048k models over the initial context length using the Hugging Face Open LLM Leaderboard~\cite{huggingfaceboard}, considering both zero-shot and few-shot evaluation paradigms. Specifically, we employ a 25-shot setting for ARC-Challenge~\cite{allenai:arc}, 10-shot configuration for HellaSwag~\cite{zellers2019hellaswag}, 5-shot regime for MMLU~\cite{mmlu}, and a zero-shot approach for TruthfulQA~\cite{lin2021truthfulqa}.

Table~\ref{tbl:commonsense} illustrates that our models perform on par with the original benchmark, which was developed for use with a more limited context window, and surpass the original Mistral on TruthfulQA by a margin of +0.5\%. While LongRoPE-LLaMA2-2048k—fine-tuned on 256k—experiences a modest decline in performance, its results largely stay within acceptable bounds across the evaluated tasks.

\subsection{Ablation Results}

\textbf{Effectiveness of the second positional interpolation}. As part of our progressive extension framework, we employ our search-based algorithm to apply a second, non-uniform positional interpolation to already fine-tuned, extended LLMs. To assess the utility of this approach, we experiment using our fine-tuned LLaMA2-256k model, incrementally extending its context length to 512k, 1024k, and 2048k with both PI and YaRN as baselines. According to Table~\ref{tbl:secondsearch}, our method of non-uniform positional interpolation preserves a stable perplexity across expansions, whereas the perplexity for PI and YaRN escalates rapidly as the context window grows.

\textbf{Recovery at shorter context lengths.} To address performance degradation at reduced context lengths, we optimize the RoPE scaling factors for LongRoPE-2048k using our search algorithm. In this process, we limit the maximum scale factors considered, thereby reducing interpolation effects at 4k and 8k context sizes. 
 Table~\ref{tbl:shortperformance} provides a perplexity comparison for LongRoPE-LLaMA2-2048k on Proof-pile at the 4k and 8k context lengths, in addition to reporting mean LLM benchmark accuracy. The data indicate that our approach leads to notable gains in performance for shorter contexts.

\textbf{Analysis on the two forms of non-uniformities}. In this section, we conduct an ablation study to investigate the individual effects of the two forms of non-uniformities on model performance. We design two experimental setups: (i) expanding LLaMA2-7B to handle 16k and 32k sequence lengths using various strategies—including PI, searching exclusively over RoPE dimension, and searching across both non-uniformities; (ii) scaling up our fine-tuned 256k-length LLaMA2 model to process up to 2048k, employing the same methodology. Perplexity measurements are obtained in the absence of further fine-tuning. According to the results presented in Table~\ref{tbl:searchdimension}, introducing non-uniformity in the RoPE dimension markedly lowers perplexity relative to the baseline PI linear interpolation. Adjusting token position non-uniformity yields noticeable gains at sequence lengths of 16k and 32k; however, its influence wanes at the extended length of 2048k, which may be attributed to the challenges inherent in extremely long-context processing. Relying solely on initial token preservation without interpolation is rendered ineffective in such scenarios, and we propose to examine this issue in subsequent research.

\section{Related Works}

Beyond position interpolation techniques, this section reviews alternative strategies present in related literature.

\textbf{Retrieval-based} approaches employ an external memory system to store information from earlier contexts, along with retrieval modules to access relevant documents during inference~\cite{longllama,wang2023augmenting,borgeaud2022improving}. Such frameworks generally require direct architectural alterations to the underlying LLM. In comparison, our method introduces only minimal changes to the positional embedding component, resulting in a more lightweight solution. Furthermore, our approach extends naturally to a broader range of long context applications, such as comprehensive document summarization and few-shot learning, not limited to retrieval tasks.

\textbf{Attention-based context window extensions}. In addition to positional embedding interpolation, certain studies have extended the input context using the inherent context window length of the original LLMs by altering the underlying attention mechanisms~\cite{han2023lminfinite, streamingllm,parallelwindow}. The central concept involves alleviating the attention explosion problem that arises from novel position indices by introducing innovative attention masks. These techniques complement the methods based on positional interpolation.

\textbf{Fine-tuning based approaches} aim to enhance the capability of pre-trained LLMs to process extended contexts by adjusting position embeddings and subsequently fine-tuning on longer sequences. Approaches such as Code LLaMA~\cite{codellama}, LLaMA2 Long~\cite{xiong2023effective}, and ScaledRoPE~\cite{liu2023scaling} employ a substantially increased RoPE base and adapt the models by fine-tuning for the desired sequence lengths. Our approach is adaptable to a wide range of target sequence lengths and is capable of handling contexts exceeding 2M tokens. In light of the intensive GPU requirements associated with fine-tuning models for very long contexts (greater than 128k), recent methods like LongLoRA~\cite{longlora} and PoSE~\cite{zhu2023pose} have been introduced to alleviate computational costs. Our approach operates independently and can complement such efficient fine-tuning strategies.

\section{Conclusion}

In this study, we introduce LongRoPE, a technique that significantly expands the context window of LLMs to a groundbreaking 2048k, all while preserving their performance on the original, shorter context ranges. By leveraging two distinct types of non-uniformity in RoPE positional embeddings, identified through a streamlined evolutionary search process, our approach yields dual advantages: it offers effective initialization for subsequent fine-tuning and allows for an eightfold increase in context length even without any additional fine-tuning. Furthermore, we devise a progressive extension methodology, utilizing LLMs fine-tuned with a 256k context length, to achieve a 2048k context window without necessitating further fine-tuning steps. Comprehensive experimental results substantiate the efficacy of LongRoPE. We anticipate that LongRoPE-2048k models will unlock a range of novel long-context applications and stimulate continued advances in this domain.

\section*{Broader Impacts} The aim of this paper is to contribute to the progression of Machine Learning as a discipline. While the outcomes of our research may influence society in various ways, we do not identify any particular issues that warrant explicit mention at this time.

	\balance

\end{document}